% !TEX root = ../thesis.tex
%
% configurations for the claude-tunnel bachelor thesis
%

% English Language support (main language stays German)
\fullenglish{no}

% Supervisors — PLACEHOLDER: replace with the actual examiners.
\firstSupervisor{Prof. Dr. Vorname Nachname}
\secondSupervisor{Prof. Dr. Vorname Nachname}

% Thesis title
\thesisTitle{Providerblinder, zensurresistenter Netzwerk-Tunnel:\\ Entwurf, Implementierung und Parameterstudie}
\thesisTitleEN{Provider-blind, Censorship-resistant Network Tunnel: Design, Implementation and Parameter Study}

% Keywords
\keywordsDE{Zero-Knowledge, providerblind, QUIC, Noise, Proof-of-Work, Zensurresistenz, NAT-Traversal}
\keywordsEN{zero-knowledge, provider-blind, QUIC, Noise, proof-of-work, censorship resistance, NAT traversal}

% Abstract
\abstractDE{Diese Arbeit entwirft, implementiert und evaluiert einen
providerblinden, zensurresistenten Netzwerk-Tunnel. Der Betreiber der
Vermittlungsebene relayt ausschliesslich Chiffretext (Ende-zu-Ende-Verschluesselung
mit dem Noise-Protokoll ueber QUIC), kennt weder Ziel-Hostnamen noch Nutzlast
und ist damit strukturell blind gegenueber dem vermittelten Verkehr. Ein
Proof-of-Work-gestuetztes Rendezvous begrenzt Missbrauch, ein direkter
Peer-to-Peer-Pfad umgeht bei Erreichbarkeit den Relay, und ein
TLS-ueber-TCP-Fallback traegt denselben Tunnel, wenn UDP blockiert ist. Eine
Parameterstudie am fertigen Produkt vermisst die Roundtrip-Latenz unter
emulierten Netzbedingungen ueber die Betriebsarten One-shot, Streaming und
Datagramm.}
\abstractEN{This thesis designs, implements and evaluates a provider-blind,
censorship-resistant network tunnel. The operator of the mesh plane relays only
ciphertext (end-to-end encryption with the Noise protocol over QUIC), knows
neither the destination hostname nor the payload and is thus structurally blind
to the traffic it forwards. A proof-of-work rendezvous limits abuse, a direct
peer-to-peer path bypasses the relay when reachable, and a TLS-over-TCP fallback
carries the same tunnel when UDP is blocked. A parameter study on the finished
product measures round-trip latency under emulated network conditions across the
one-shot, streaming and datagram operating modes.}

% Author — PLACEHOLDER: replace with the student's full name.
\thesisAuthor{Vorname Nachname}

% Submission date — PLACEHOLDER
\submissionDate{TT. Monat 2026}

% NDA / public
\NDA{no}

% Cover style
\Cover{CD2017}

% Lists
\ListOfFigures{yes}
\ListOfTables{yes}
\ListOfAccronyms{yes}
\ListOfSymbols{no}
\Glossary{yes}

% Study course
\studycourse{INF_AI}   % Angewandte Informatik
